\documentclass[12pt]{article}
\usepackage{amsmath, amssymb, amsthm}
\usepackage{parskip}
\usepackage{geometry}
\usepackage{enumitem}
\usepackage{fancyhdr}

% Set margins to match a standard homework document
\geometry{a4paper, margin=1in}

\title{Recitation 4}
\parskip=1ex
\author{Math 311 Fall 2026 \\ Rutgers University}
\date{}

\begin{document}

\maketitle

\textbf{Exercise 1:}
Let \(x_1 > 0\) and define
\[
    x_{n+1} = \frac{1}{2}(x_n + \frac{3}{x_n}) \qquad \forall n \in \mathbb{N}.
\]
\begin{enumerate}[label=(\alph*)]
    \item
          Since \(x_1 > 0\), suppose \(x_n > 0\) and so
          \[
              x_{n+1} = \frac{1}{2}(x_n + \frac{3}{x_n}) > 0.
          \]
          Thus, $x_n > 0$ for all \(n \in \mathbb{N}\) by induction.
          Then
          \[
              \begin{split}
                  x_{n+1} - \sqrt{3} & = \frac{1}{2}\left(x_n + \frac{3}{x_n}\right)- \sqrt{3} \\
                                     & = \frac{x_n^2 + 3 - 2\sqrt{3}x_n}{2x_n}                 \\
                                     & = \frac{\left(x_n - \sqrt{3}\right)^2}{2x_n}.
              \end{split}
          \]
          Since the numerator is squared, it is non-negative,
          and we already showed that the \(x_n > 0\), so \(2x_n > 0\) as well.
          Thus \(x_{n+1} - \sqrt{3} \geq  0 \), and so
          \[
              x_{n} \geq \sqrt{3}.
          \]
          for every \(n \geq 2\).
    \item Consider
          \[
              \begin{split}
                  x_{n+1} - x_n & = \frac{1}{2}\left(x_n + \frac{3}{x_n}\right) - x_n \\
                                & = \frac{3}{2x_n} - \frac{1}{2} x_n                  \\
                                & = \frac{3-x_n^2}{2x_n}.
              \end{split}
          \]
          Since $x_n \geq \sqrt{3}$ for every \(n \geq 2\), we have
          \[
              3-x_n^2 \leq 0 \quad \text{and} \quad 2x_n > 0.
          \]
          Then
          \[
              x_{n+1} \leq x_n.
          \]
          Hence the sequence is monotonically decreasing and bounded below, so by the Monotone Convergence Theorem,
          \(\{x_n\}\) converges to some limit \(L\). Taking the limits of the recurrence we have
          \[
              L = \frac{1}{2}\left(L + \frac{3}{L}\right).
          \]
          Which simplifies to
          \[
              L^2 = 3.
          \]
          Therefore,
          \[
              \{x_n\} \to \sqrt{3}.
          \]












\end{enumerate}


\end{document}

